% Round 16, global_height_audit. New bridge; independent audit requested.
% Historical sources are inputs with their frozen scope, not modified here.
\section{From a vanishing actual obstruction to reciprocal Green pairings}
\label{r16h:section}

\paragraph{Status and conventions.}
The point-slice and height theorems below have passed independent review
at root and in the model-obstruction branch, relative to their stated
frozen inputs. The review is recorded in
\path{round16/audits/obstruction_zero_height_model_review.md}; the
finite-flat base-change step was additionally checked by the trace
branch. These are internal mathematical reviews, not external peer
review or formal verification. This proof uses the Round 16
full finite actual trace reconstructed from the final Round 11 interfaces;
it does not presume the historical Round 6 trace candidate. The latter
reconstruction has been independently checked by this author in
\texttt{global\_height/full\_finite\_trace\_independent\_audit.md}.
The new result below is a proper-slice residue theorem for the concrete
Round 11 supported-complex Green pair. In particular it uses
$\mathfrak o=0$, not an unproved converse $G\Rightarrow\mathfrak o=0$.

Fix a finite extension $K/\mathbf Q_p$ with its original absolute value.
All forms are Mihatsch actual forms in the ordinary topology, and
$D=d'd''$. Integration uses the native $K$-weights. A rational cycle is
always written with its final nonzero prime coefficients before choosing
its supporting complexes.

\subsection{The finite-flat interface, including ramification}

\begin{lemma}[Finite-flat base change for an arbitrary normal testing map]
\label{r16h:finite-base-change}
Let $q:V\to U$ be finite flat of degree $d$, with $U,V$ smooth integral
$K$-varieties and $U$ geometrically integral. Let $a:M\to U$ be a
morphism with $M$ normal integral. Write the fundamental cycle of the
finite flat fiber product as
\[
 [V\times_U M]=\sum_j e_j[W_j],
 \qquad N_j=W_j^\nu,\quad q_j:N_j\to M,\quad \rho_j:N_j\to V,
\]
where $e_j$ is the generic Artinian length. For every actual form $\omega$
on $V^{\rm an}$,
\[
 a^*\operatorname{Tr}_q\omega
   =\sum_j e_j\operatorname{Tr}_{q_j}(\rho_j^*\omega).
 \tag{H16.1}\label{r16h:base-change}
\]
This is an equality of ordinary actual forms, in every bidegree. It is
local on ordinary opens and retains all native field degrees.
\end{lemma}

\begin{proof}
We first record why taking a product commutes with trace here, without
assuming a general current base-change theorem. In the explicit Round 16
construction, a source presentation $\omega=t_f^\star\gamma$ is encoded by
the independent-parameter kernel
$\Gamma_f=\prod_iD\max(v(f_i)-v(z_i),0)$. The trace is the supplied
product-moment integral of
$\chi(t)\gamma(t)\operatorname{Tr}^{\rm curv}_{q\times\mathrm{id}}\Gamma_f$.
After taking a product with a smooth variety $M_0$, pull back exactly the
same source units, Newton coefficients, descended sorted potentials and
parameter cutoff. The pure-curvature orbit identity is unchanged.
To check its trace pairing, pass to a normal Galois closure and use the
native finite degree formula on its product with $M_0$. These products
are normal because $M_0$ is smooth. If they are disconnected, sum over
their disjoint normal integral components; their degrees over each
target component sum to the original degree. The Galois action permutes
these components, so the same orbit-average proof applies. Thus the
pulled pure-curvature output is the product-map pure-curvature trace.
The supplied-presentation base-change identity FT16.3 then gives
\[
 \operatorname{Tr}_{q\times\mathrm{id}_{M_0}}(\mathrm{pr}_V^*\omega)
             =\mathrm{pr}_U^*(\operatorname{Tr}_q\omega).
 \tag{H16.2}\label{r16h:product-trace}
\]
This argument uses the explicit trace formula and actual uniqueness,
not a proposed base-change law for arbitrary currents.

For smooth $M$, factor $a$ as its closed graph in $U\times M$ followed
by the projection. Apply the ramified closed-slice proof in
\texttt{round10/trace/ramified\_form\_base\_change.tex},
Theorem G10.T2, to $q\times\mathrm{id}_M$ and this graph, and use
\eqref{r16h:product-trace}. Its formerly conditional full-trace input is
now supplied by Round 16. For clarity, the essential calculation of that
proof is the following. Locally a smooth closed slice is a regular
sequence $f_1,\ldots,f_c$. Mihatsch's complete-intersection kernel
\[
 g=(-1)^{c-1}\phi(D\phi)^{c-1},\qquad \phi=\min_i v(f_i),
\]
has the strong all-actual-test residue identity; flat pullback preserves
the regular sequence, including its scheme multiplicities. For a compact
actual form $\beta$ of the slice's top degree, $D\beta$ has support away
from that fiber product. A compact cutoff of $g$ equal to one over the
image of this support is therefore an allowed trace test. Trace
adjunction and commutation with $D$ give
\[
 \int_{\text{slice}}\operatorname{Tr}\beta
                   =\int_{\text{fiber product}}\beta.
\]
Arbitrary actual slice tests lift locally by lifting their unit moments
and keeping their original Euclidean delta coefficient. Multiplication
by these lifted tests gives the form-valued assertion. Decomposing the
fiber-product fundamental cycle and normalizing produces exactly the
generic lengths $e_j$, not coefficient one. No trace on a nonreduced
or nonnormal target is introduced.

For normal $M$, the right side of \eqref{r16h:base-change} exists on all
of $M$: finite flatness makes every $W_j$ dominate $M$, and normalization
is finite. Compare on the dense smooth algebraic locus of $M$ by the
preceding argument. Actual forms on a pure-dimensional variety are
determined on a dense algebraic open. Indeed, a difference supported
on a proper algebraic closed subset, wedged with any complementary
smooth compact test, is a top actual form whose support avoids that
subset by Mihatsch Corollary 4.9. Its current vanishes, and Proposition
4.11 gives actual vanishing. This proves the equality everywhere.
\end{proof}

\subsection{A supported bivariant calculation}

\begin{lemma}[The point residue of a supported perfect complex]
\label{r16h:algebraic-residue}
Let $X/K$ be smooth of dimension $n$, let $Z\subset X$ have pure
codimension $r$, and let $C$ be a perfect complex exact off $Z$. Put
\[
 c=\operatorname{ch}_r^Z(C),\qquad z_C=c\cap[X]\in\operatorname{CH}_{n-r}(Z)_{\mathbf Q}.
\]
For an integral $r$-dimensional variety $a:M\to X$ such that
$S=M\times_X Z$ is finite, one has
\[
 \operatorname{ch}_r^S(La^*C)\cap[M]
       =\Delta_X^!([M]\times z_C)
       =:I_a(z_C)\quad\text{in }\operatorname{CH}_0(S)_{\mathbf Q}.
 \tag{H16.3}\label{r16h:refined-residue}
\]
Because $S$ is finite, this is equality of the specified point
coefficients, including nonreduced and residue-field contributions.
\end{lemma}

\begin{proof}
Localized Chern characters are bivariant classes and commute with
arbitrary base change (Stacks 0FB3). The graph of $a$ is a regular
immersion in $X\times_K M$, since it is a section of the smooth
projection to $M$. Bivariant classes commute with its refined Gysin map
(Stacks 0FBP). Pull $c$ to $X\times M$; its action on $[X]\times[M]$
is $z_C\times[M]$, by flat exterior-product compatibility. Applying
the graph Gysin and commuting it past $c$ gives the left side of
\eqref{r16h:refined-residue}; the graph Gysin of $[X]\times[M]$ is
$[M]$. Its action on $z_C\times[M]$ is precisely the refined diagonal
intersection on the right. This is the smooth-orientation calculation
of Stacks 0FBR, Proposition 42.60.2; the associated intersection is
defined in Stacks 0FC0. Finally $\operatorname{CH}_0$ of a finite
scheme is the free group on its closed points, so no rational
equivalence can hide a changed point coefficient.
\end{proof}

\subsection{The new curved proper-slice theorem}

Use the precise Round 11 supported-model-complex setup:
$\mathcal X/R$ is integral flat projective with smooth geometrically
integral generic fiber $X$, $C$ is a bounded vector-bundle complex on
$\mathcal X$ exact off its full model cohomology support, and its generic
support $Z$ is pure of codimension $r$. Its frozen Green pair is
\[
 \Omega_C=\sum_i(-1)^i\operatorname{ch}_r(\overline C_i),
 \qquad G_C\in B^{r-1,r-1}((X\setminus Z)^{\rm an}),
 \qquad D[G_C]+\delta_{z_C}=[\Omega_C].
\]
The last identity already holds on every compact actual test on $X$.
The source is \texttt{round11/perfect\_complex/supported\_complex\_green.tex},
PC2--PC10, together with the final minimal-trace audit. We retain its
particular construction, not an arbitrary solution of the current equation.

\begin{theorem}[All-actual finite-intersection residue identity]
\label{r16h:curved-slice}
Let $a:M\to X$ be any morphism from a normal integral quasiprojective
$K$-variety of dimension $r$, with finite $S=a^{-1}Z$. The particular
Round 11 form $G_C$ satisfies
\[
 \int_{M\setminus S}a^*G_C\wedge D\eta
       +\int_{I_a(z_C)}\eta
       =\int_M a^*\Omega_C\wedge\eta
 \tag{H16.4}\label{r16h:curved-slice-eq}
\]
for every compact actual scalar test $\eta\in B_c^{0,0}(M^{\rm an})$.
The form $a^*G_C$ is ordinary actual on exactly $M\setminus S$.
\end{theorem}

\begin{proof}
The relative diagonalization PC2 supplies a proper generically finite
$f:T\to X$, of degree $d$, with $T$ smooth projective and $f$ finite
over $U=X\setminus Z$, and a supported perfect-$K_0$ identity
\[
 [f^*C]=\sum_\alpha\epsilon_\alpha[L_\alpha\xrightarrow{s_\alpha}M_\alpha].
\]
Every $s_\alpha$ is an isomorphism off $f^{-1}Z$. The finite restriction
is flat because source and target are smooth of equal dimension. Put
\[
 B_{r-1}(l,m)=\frac1{r!}\sum_{b=0}^{r-1}m^{r-1-b}l^b,
 \quad g_\alpha=-\log\|s_\alpha\|,\quad
 G_T=\sum_\alpha\epsilon_\alpha g_\alpha B_{r-1}(l_\alpha,m_\alpha).
\]
Then $G_C=d^{-1}\operatorname{Tr}_{f|_{f^{-1}U}}G_T$. The minimal
trace and the full trace agree on this scalar-curvature input, since
both are actual and have the same all-actual-test adjunction.

In $T\times_X M$, take the normalizations $N_j$ of the reduced
components dominating $M$, with their generic Artinian lengths $e_j$.
Write $b_j:N_j\to M$ and $\rho_j:N_j\to T$. These are proper over
$M$, are finite over $M\setminus S$, and
\[
                  \sum_j e_j(b_j)_*[N_j]=d[M].
 \tag{H16.5}\label{r16h:source-degrees}
\]
Extra components lying over $S$ are not included; they contribute zero
to this displayed pushforward. The pulled line maps on each $N_j$ are
nonzero, since their generic points lie over $U$. Their Cartier
divisors $D_{\alpha j}$ lie over the finite set $S$.

On each normal $N_j$, apply scalar strong Poincare--Lelong to the
pulled line maps with the compact actual test
$\rho_j^*B_{r-1}(l_\alpha,m_\alpha)b_j^*\eta$. Compactness follows
from properness of $b_j$. The Chern factors are separately closed.
The strict curvature identity PC9 pulls back as an actual identity.
Sum with weights $e_j\epsilon_\alpha$ to obtain
\begin{align*}
 &\sum_j e_j\int_{N_j\setminus b_j^{-1}S}
                    \rho_j^*G_T\wedge b_j^*D\eta\\
 &\quad+\sum_{j,\alpha}e_j\epsilon_\alpha
       \int_{D_{\alpha j}}\rho_j^*B_{r-1}(l_\alpha,m_\alpha)b_j^*\eta
       =\sum_j e_j\int_{N_j}b_j^*(a^*\Omega_C\wedge\eta).
 \tag{H16.6}\label{r16h:upstairs-slice}
\end{align*}
Scalar strong Poincare--Lelong applies to normal integral sources; the
argument has not asserted that these $N_j$ are smooth.

The actual $(1,1)$ form $D\eta$ has support disjoint from the finite
algebraic set $S$ by Corollary 4.9. Its pulled support is compact in
the finite locus. Lemma~\ref{r16h:finite-base-change} and finite trace
adjunction identify the first line with
$d\int_{M\setminus S}a^*G_C\wedge D\eta$.
The top actual form $a^*\Omega_C\wedge\eta$ is likewise supported
away from $S$. The native degree formula and
\eqref{r16h:source-degrees} identify the right side with
$d\int_Ma^*\Omega_C\wedge\eta$.

Each component of $D_{\alpha j}$ is projective over its image, a
closed point of $S$. Its Chern polynomial has top degree $r-1$.
Thus the frozen Chern-degree lemma PC6, now with zero-dimensional
target, identifies its analytic integral with the algebraic Chern
degree times the value of $\eta$, with the native residue degree.
The Cartier localized-character identity PC4 and the proper
bivariant projection formula therefore identify the middle sum with
the integral of the zero-cycle
\begin{align*}
 &\sum_j e_j(b_j)_*\left(
           \operatorname{ch}^{b_j^{-1}S}_r(Lb_j^*La^*C)\cap[N_j]\right)\\
 &\hspace{25mm}=\operatorname{ch}_r^S(La^*C)\cap
                    \sum_j e_j(b_j)_*[N_j]
       =d I_a(z_C).
\end{align*}
Here the supported line-complex $K_0$ identity was pulled to $N_j$;
localized characters commute with this arbitrary pullback. The final
equality is Lemma~\ref{r16h:algebraic-residue}. Everything takes place
in $\operatorname{CH}_0(S)$, hence gives the exact point coefficients.
Divide \eqref{r16h:upstairs-slice} by $d$.

No general proper actual pushforward appears: the only operations are
finite actual traces away from $S$, scalar Poincare--Lelong on normal
sources, and scalar top-Chern integration on proper fibers.
\end{proof}

\begin{corollary}[The obstruction-zero sector has the full point-slice property]
\label{r16h:P}
Let $Z$ be a specified rational codimension-$r$ cycle on a smooth
projective geometrically integral $X/K$, and assume
$\mathfrak o_X^r([Z])=0$. There is an ordinary actual strong Green form
$\gamma_Z$ on exactly $X^{\rm an}\setminus|Z|^{\rm an}$ which, for
every normal integral quasiprojective $r$-fold $a:M\to X$ with finite
$S=a^{-1}|Z|$, satisfies
\[
             D[a^*\gamma_Z]=-\delta_{I_a(Z)}
 \tag{H16.7}\label{r16h:P-eq}
\]
on all compact actual scalar tests. In particular the frozen
$P_{\rm chain}$ condition holds. The same chosen form works for all
such maps $a$.
\end{corollary}

\begin{proof}
Choose the Round 11 supported complexes component by component, with
the original final support, obtaining $(G_Z,\Omega_Z)$ and $z_C=Z$.
The definition and Chow-class independence of the frozen obstruction
give a global actual $\Psi$ with $D\Psi=\Omega_Z$. Set
$\gamma_Z=G_Z-\Psi|_{X\setminus|Z|}$. The already proved Round 11
Green equation gives strong $G$ on all actual tests, with the exact
support. Pullback of the globally defined actual primitive commutes
with $D$. Subtract its equation from
Theorem~\ref{r16h:curved-slice}, component by component.
This is a concrete residue proof; it does not pull back a distributional
Green equation through an intersection.
\end{proof}

\subsection{Finite constant extensions of the concrete Green choice}

\begin{proposition}[Absolute normalization for obstruction-zero choices]
\label{r16h:constant-extension}
For a finite extension $L/K$, a Green form chosen as in
Corollary~\ref{r16h:P} pulls back to a strong Green form for $Z_L$.
If $Z,W$ have complementary codimensions, disjoint supports, and both
local obstruction classes vanish, then
\[
                   \lambda_L(Z_L,W_L)=\lambda_K(Z,W).
 \tag{H16.8}\label{r16h:local-absolute}
\]
\end{proposition}

\begin{proof}
Base change the supported model complex, relative line decomposition,
and its model metrics to $R_L$. Generic sources remain smooth, possibly
with several disjoint components, whose total degree is $d$. Lemma
\ref{r16h:finite-base-change}, applied on the finite locus to
$X_L\to X$ viewed as a $K$-morphism, says that the pullback of $G_C$
is exactly the base-changed line-sum trace divided by $d$. The proof
PC4--PC10 therefore gives its all-actual Green equation with cycle
$Z_L$ and curvature $\Omega_{C,L}$. The actual primitive pulls back
with $D\Psi_L=\Omega_{C,L}$, so subtraction gives strong $G$.

The analytic space of an $L$-variety viewed over $K$ has the same
ordinary actual forms: its analytic units and genuine moment
presentations are unchanged. Its native $K$-integrals are $[L:K]$
times its native $L$-integrals. This is the constant-field weight
normalization in Mihatsch Lemma 3.16 and Section 4.3. Dividing the
base-changed equations by this common factor preserves every actual
test and gives the asserted equation over $L$.

The native finite degree formula for $X_L\to X$, followed by that
same factor cancellation on $W_L$, gives
$\int_{W_L/L}\gamma_{Z,L}=\int_{W/K}\gamma_Z$. Disjoint-polar
Stokes makes both local numbers independent of the chosen Green
forms because both complementary cycles have strong solutions. This
proves \eqref{r16h:local-absolute}.
\end{proof}

\subsection{A canonical real height on the local obstruction-zero domain}

Let $X/F$ be smooth projective geometrically integral of dimension $n$
over a number field. Geometric homological triviality is with rational
coefficients. Fix the usual normalized place absolute values and put
$n_v=[F_v:\mathbf Q_v]/[F:\mathbf Q]$. Define the Chow subspace
\[
 \mathcal H^r(X)=\{\alpha\in\operatorname{CH}^r(X)_{\mathbf Q,\rm hom}:
       \mathfrak o_{X_{F_v}}^r(\alpha_{F_v})=0
                    \text{ for every finite place }v\}.
 \tag{H16.9}\label{r16h:height-domain}
\]
Outside a fixed finite set of bad-reduction places, the local
obstruction vanishes for every such class by the frozen Round 11
smooth-projective-good-model theorem. Thus this definition imposes
only finitely many possibly nonzero local obstruction maps.

\begin{theorem}[Obstruction-zero analytic Beilinson--Bloch height]
\label{r16h:global-height}
For $p,q\geq1$ with $p+q=n+1$, there is a canonical real-valued
$\mathbf Q$-bilinear pairing
\[
 h_X^{\rm an}:\mathcal H^p(X)\times\mathcal H^q(X)\longrightarrow\mathbf R.
 \tag{H16.10}\label{r16h:height}
\]
For disjoint representatives $Z,W$, it is the finite sum
\[
 h_X^{\rm an}([Z],[W])
         =\sum_v n_v\lambda_v(Z_v,W_v),\qquad
 \lambda_v(Z_v,W_v)=\int_{W_v^{\rm an}}\gamma_{Z_v}\quad(v<\infty),
 \tag{H16.11}\label{r16h:height-sum}
\]
with the classical zero-curvature complex Green normalization at
infinite places. It is independent of all Green choices and moving
representatives, symmetric under interchange of the two factors, and
invariant under finite number-field extension. It is the adelic degree
of the canonical Green metric on the corresponding restriction of
Bloch's algebraic biextension, with the dual/sign convention fixed
by the scalar Green function $-\log|f|$.
\end{theorem}

\begin{proof}
\emph{Local definition and symmetry, with its current hypotheses.}
Every specified representative of a class in $\mathcal H^r$ has the
strong Green form of Corollary~\ref{r16h:P}. Here the precise frozen
input is Lemma \texttt{lem:extension}, Proposition
\texttt{prop:stokes}, and Corollary \texttt{cor:independence} in
\texttt{history\_01\_07/analytic\_bb\_input/}\allowbreak
\texttt{Lecture\_F\_Analytic\_Conjectures\_Beilinson\_Bloch\_2026-11-04.tex}.
Their hypotheses hold for every strong Green choice in the question,
not only the constructed one. Indeed, if $S$ has codimension at least
$r$, every compact complementary actual test for a form of degree
$(r-1,r-1)$ has degree $(n-r+1,n-r+1)$ and support disjoint from
$S$ by Corollary 4.9. Its integral with any off-$S$ actual Green
form is consequently a finite integral on a compact subset away
from $S$. It defines the canonical algebraic-dual actual current;
on original smooth tests it is continuous by Mihatsch Lemma 5.1.
No extra local integrability hypothesis at $S$ is needed.

For disjoint polar supports $S,T$, the opens $X\setminus S$ and
$X\setminus T$ cover $X$. On either open at least one Green current
is a genuine actual form, so the module action of actual forms on
the algebraic dual of compact actual tests defines the mixed
currents there, and their ordinary form products agree on the overlap.
If $g=[\gamma_Z]$ and $h=[\gamma_W]$, the identity
\[
 (Dg)h-g(Dh)=d'((d''g)h)+d''(g(d'h))
\]
holds by the Leibniz rule on both opens. Its right side vanishes
on the compact test $1$ because $X$ is proper and boundaryless.
Substitution of the two strong equations gives
\[
                    \int_W\gamma_Z=\int_Z\gamma_W.
\]
It also shows independence of either Green choice: subtract two
strong solutions in one factor and apply the same identity with
zero residue in that factor. This uses all actual tests, not a
model-generated subspace. Additivity follows on common disjoint
position, keeping the union of polar supports when cycle components
cancel. The archimedean zero-curvature construction has the same
properties.

\emph{Almost all local terms vanish.}
Spread a fixed disjoint pair and $X$ to a smooth projective model
over $\mathcal O_F[S^{-1}]$. After enlarging $S$, their flat closures
are disjoint: the intersection is proper over this one-dimensional
base with empty generic fiber, hence has finite image. On a good
local model choose finite vector-bundle resolutions of those
closures, on the fixed regular projective model itself. Their full
cohomology supports are the respective closures. The Round 11
good-model theorem gives zero model curvature for the homologically
trivial generic cycle and a strong Green form literally zero on
every tube disjoint from its closure. The other cycle lies in those
tubes, so its integral is zero. Green-choice independence transfers
this value to the originally chosen form. Thus
\eqref{r16h:height-sum} is a finite sum.

\emph{The exact principal transport.}
Suppose $Z'-Z=\partial B$, with $Z,Z'$ disjoint from $W$. The good
higher-Chow moving complex supplies such a one-chain $B$ over $F$,
with both faces separately
avoiding $W$ and finite proper intersection with $W$ along the
normalized chain. For this calculation first clear denominators in
\emph{both} $Z,Z',B$ and $W$. Local rational homogeneity has already
been proved, so one divides by the product of the two chosen
multipliers at the end. All point exponents in the norm computation
that follows are therefore integers. For one integral chain component, normalize its
projective closure in $X\times\mathbf P^1$. Remove the fiber of
the parameter $f$ at $-1$, obtaining a normal $q$-fold $a:M\to X$.
Its inverse image $S=a^{-1}|W|$ is finite, $f$ is a unit near $S$,
and $a_*\operatorname{Div}_M(f)=\partial B$. The precise moving
input is Gorchinskiy's Lemma 2.1: the subcomplex meeting the fixed
components of $|W|$ and all simplex faces properly is
quasi-isomorphic to the full higher-Chow complex. Since $Z'-Z$ is
already in its degree-zero part and is a boundary in the full
complex, injectivity on degree-zero homology produces the required
chain over the original field $F$. The faces each have expected
intersection dimension $-1$ with $W$, hence separately avoid it.
Use $\Delta^1=\{t_0+t_1=1\}$, $f=t_0/t_1$, and
$\partial B=B\cdot\{t_0=0\}-B\cdot\{t_1=0\}$; then the omitted
fiber is $f=-1$ and the divisor has the displayed sign.
The scalar Green function $-\log|f|$ is zero on the ordinary open
$|f+1|<1$ surrounding the removed fiber, so its current and
derivatives have compact support on $M$. Apply scalar
Poincare--Lelong and compact-support disjoint-polar Stokes to this
function and the actual pullback Green form for $W$ supplied by
\eqref{r16h:P-eq}. If
$I_a(W)=\sum_{x\in S}m_x[x]$, the result is
\[
 \lambda_v(\partial B,W)
       =-\log|a(B,W)|_v,\qquad
 |a(B,W)|_v=
 \left|\prod_{x\in S}\operatorname{Norm}_{\kappa(x)/F}(f(x))^{m_x}\right|_v.
 \tag{H16.12}\label{r16h:principal}
\]
Here $a(B,W)\in F^\times$ is the \emph{algebraic} higher-Chow
transport element in the chosen biextension convention. With the
standard simplex symbol written using $-f$, its point formula uses
$\operatorname{Norm}(-f(x))$ instead. The two displayed point products
differ only by a sign, and $|-1|_v=1$ at every place. Any omitted
point evaluation at $f=-1$ likewise has absolute value one.
Thus the asserted equality of logarithms does not require an
incorrect identification of differently signed field elements.
For a sum of chain components multiply their algebraic norm
evaluations. This is exactly the degree-one higher-Chow product
and pushforward into $\operatorname{CH}^1(F,1)=F^\times$;
Gorchinskiy's Proposition 4.2 identifies its biextension with
Bloch's. No proper-intersection assertion for the projected
$K_1$-chain carrier is substituted for the good higher-Chow chain.
The exact residue coefficients of Theorem~\ref{r16h:curved-slice}
are what makes this the absolute value of the prescribed algebraic
norm, rather than an unspecified local constant. At completion, finite field tensors,
norms and refined Gysin coefficients base change with their full
degrees, so the right side is one global element evaluated at $v$.
The established complex Green construction has the same principal
law at infinity.

\emph{Chow descent.}
For $Z,Z'$ representing the same class and avoiding $W$, the moving
complex gives the good chain just used. Summing
\eqref{r16h:principal} with the weights $n_v$ gives zero by the
number-field product formula. Thus the global value is unchanged
by rational equivalence in the first argument; symmetry gives the
second argument. To compare two disjoint pairs $(Z_1,W_1)$ and
$(Z_2,W_2)$, choose a representative $W_0$ avoiding
$|Z_1|\cup|Z_2|$ and compare successively
$(Z_1,W_1),(Z_1,W_0),(Z_2,W_0),(Z_2,W_2)$. This proves independence
of simultaneous moving. Moving finitely many cycles into common
disjoint position and using local additivity proves rational
bilinearity. Symmetry was already established locally.

\emph{Absolute base field.}
The obstruction is functorial under finite constant extension, so
these domains map into the corresponding domains after extension.
Proposition~\ref{r16h:constant-extension} gives equality of each
local value after such extension, with the classical equality at
infinity. For $F'/F$, the identity
\[
 \sum_{w\mid v}\frac{[F'_w:\mathbf Q_v]}{[F':\mathbf Q]}
                        =\frac{[F_v:\mathbf Q_v]}{[F:\mathbf Q]}
\]
then proves invariance of the global sum.

\emph{Identification of the algebraic object.}
Bloch's biextension has the representative frame transition
$s_{Z',W}=a(B,W)s_{Z,W}$, with the dual convention corresponding
to $-\log|f|$. Define $\|s_{Z,W}\|_v=\exp(-\lambda_v(Z,W))$.
Equation \eqref{r16h:principal} is exactly the change-of-frame norm
law; local bilinearity makes both biextension laws isometries.
The algebraic reciprocity and closed-chain relations give coherence.
Almost-all-place vanishing makes a global representative frame have
norm one outside a finite set. Its adelic degree is
\eqref{r16h:height-sum}; another global frame has the same degree
by the product formula. Integral homological torsion is killed
after multiplication and the rational version is obtained by
bilinearity, not by claiming an ordinary line is a fractional
tensor power. The moving and biextension inputs are precisely
Gorchinskiy, \emph{Notes on the biextension of Chow groups}, Lemma
2.1, Section 4.1 and Proposition 4.2, as in the frozen initial
Lecture F.
\end{proof}

\paragraph{What this theorem does not assume or prove.}
It does not assume $P_{\rm chain}$, tube vanishing or scalar-extension
compatibility as extra conjectures on this domain: these were just
proved for the concrete supported-complex choices. It does not show
that $\mathcal H^r(X)$ is the complete homologically trivial Chow
group in general. Nor does it establish familywise continuity of a
biextension metric, nondegeneracy, positivity, a height conjecture,
or, by itself, comparison with a particular admissible
arithmetic-model metric. The separate independently reviewed
reconstruction \texttt{r16mm:local}--\texttt{r16mm:global} now gives
that comparison when regular numerically admissible model lifts
are supplied. It does not supply those lifts for every class in
$\mathcal H$. No historical threefold classical-height claim is
used in this proof.

\begin{corollary}[Domains supplied by the activated obstruction theorems]
\label{r16h:domains}
Relative to the internally reviewed Rounds 12, 15 and 16 inputs, the pairing above is
unconditional on the following domains.
\begin{enumerate}
\item For every smooth projective $n$-fold with $n\geq2$, the new
      one-cycle theorem \texttt{r16G:one-cycles} gives
      $\mathcal H^{n-1}=\operatorname{CH}^{n-1}_{\mathbf Q,\rm hom}$.
      Consequently the complete one-cycle factor pairs with
      $\mathcal H^2$ in the complementary codimension. This does not
      make $\mathcal H^2$ the complete group in arbitrary dimension.
\item For every smooth projective threefold over a number field, it
      is defined on the complete pair
      $\operatorname{CH}^2_{\mathbf Q,\rm hom}\times
       \operatorname{CH}^2_{\mathbf Q,\rm hom}$, by Round 15 G1.
\item For a supplied projective-bundle tower over a smooth projective
      geometrically integral base of dimension at most three, it is
      defined on both complete complementary homologically trivial
      Chow groups in every admissible pair of codimensions, by
      Round 15 G3 and its actual obstruction primitives. The same is
      now true for a smooth projective variety admitting, after a
      finite number-field extension, an everywhere-defined projective
      dominant morphism from such a tower, by
      \texttt{r16G:towers}. This assertion does not replace that
      morphism by an arbitrary rational map.
\item On every geometrically uniruled fourfold, Round 15 G2 supplies
      the complete codimension-two factor and Round 16's one-cycle
      theorem supplies the complete codimension-three factor. Thus
      the pairing is defined on
      $\operatorname{CH}^2_{\mathbf Q,\rm hom}\times
       \operatorname{CH}^3_{\mathbf Q,\rm hom}$ in full. The two
      endpoint codimensions are also complete by the frozen divisor
      and degree-zero zero-cycle theorems, so all admissible
      complementary pairs on this fourfold are covered.
\item For every geometrically unirational variety of dimension at
      most five, both complete complementary homologically trivial
      Chow groups are covered, by the independently reviewed
      theorem \texttt{r16cov:unirational}. In particular every
      geometrically unirational fivefold has the full middle
      self-pairing
      $\operatorname{CH}^3_{\mathbf Q,\rm hom}\times
       \operatorname{CH}^3_{\mathbf Q,\rm hom}\to\mathbf R$.
      More generally, the same full conclusion holds for targets
      of the rational domination of dimension-at-most-five sources
      satisfying the original-field full-kernel hypothesis of
      \texttt{r16cov:rational-low-dimension}. In higher dimensions,
      one applies \texttt{r16cov:rational-outer} separately to
      both complementary codimensions; a theorem for one factor
      alone does not supply the other factor.
\item In every dimension, geometrically algebraically trivial classes
      lie in $\mathcal H^r$ in all codimensions, by that same Round 12
      theorem. Either factor can also be enlarged by any independently
      proved local obstruction-zero sector in its own codimension.
\item The complete complementary groups are covered for every
      variety admitting, after a finite number-field extension, the
      supplied Chow-motive factorization through Tate twists of
      abelian varieties in \texttt{r16ab:motivic}. In particular this
      includes all abelian varieties and arbitrary finite products
      of smooth projective geometrically integral curves, in every
      dimension and codimension. The factorization is an identity
      of rational Chow correspondences, not merely a decomposition
      in geometric cohomology. Its local original-field conditions
      are checked in \texttt{r16ab:local-G}.
\end{enumerate}
These are statements about the above canonical real analytic
biextension pairing; no bibliographic priority claim is made.
\end{corollary}

The new one-cycle and dominant-tower inputs, including their exact
native-field conditions, are in
\path{round16/model_obstruction/one_cycles_and_dominant_activation.tex}.
The newly added rational-domination consequences are in
\path{round16/root/proper_covariance_and_rational_domination.tex};
their independent review is
\path{round16/audits/proper_covariance_rational_domination_model_review.md}.
The abelian and Chow-motivic inputs are in
\path{round16/model_obstruction/local_abelian_and_chow_motivic_activation.tex},
independently reviewed by the trace and global-height branches.

\begin{corollary}[Functoriality of the proved global height domains]
\label{r16h:domain-functoriality}
The groups $\mathcal H^r$ are stable under algebraic pullback,
proper pushforward between smooth projective varieties, and
algebraic correspondences, with the usual codimension shifts.
Every finite sum of correspondence images of known complete
homological obstruction-zero sectors is therefore a proved
height-domain class in its target codimension.
\end{corollary}

\begin{proof}
Algebraic pullback preserves geometric homological triviality and
commutes with every local obstruction map. At each finite place,
\texttt{r16cov:proper-push} and
\texttt{r16cov:correspondences} give the corresponding stability
of the local obstruction kernel. Algebraic proper push,
intersection and finite field base change commute, so these are
the completions of the original number-field operations.
Taking the intersection over the finite places proves the claim.
The statement concerns membership of both arguments in the
height domain; it does not claim a complete homological group
unless such a source decomposition has been supplied.
\end{proof}

\paragraph{Primary and internal sources.}
The primary algebraic sources checked for the new residue step are
\url{https://stacks.math.columbia.edu/tag/0FB3},
\url{https://stacks.math.columbia.edu/tag/0FBP},
\url{https://stacks.math.columbia.edu/tag/0FBR}, and
\url{https://stacks.math.columbia.edu/tag/0FC0}.
The all-degree Cartier localized polynomial remains the exact frozen
PC4 input, Fulton, \emph{Intersection Theory}, Corollary 18.1.2;
a leading-term localized statement alone would not suffice.
Analytic inputs are Mihatsch's actual pullback, strong scalar and
complete-intersection Poincare--Lelong, support avoidance and current
injectivity, as itemized in the cited R10/R11 proofs. Internal inputs
are final R11 minimal trace and supported complexes, the R16 full
trace, the frozen initial disjoint-polar Stokes/moving construction,
and the specific R12/R15 obstruction-vanishing theorems used above.
The moving and biextension source was independently re-opened as
Gorchinskiy, arXiv:0802.1437v2 (7 March 2009),
\url{https://arxiv.org/pdf/0802.1437}, especially Lemmas 2.1, 2.2
and Proposition 4.2. The complex Green existence and logarithmic
regularity input is Gillet--Soul\'e, \emph{Arithmetic intersection
theory}, Sections 1.2--1.3,
\url{https://www.numdam.org/article/PMIHES_1990__72__93_0.pdf},
with its scalar normalization rescaled to $-\log|f|$ as in the
frozen initial Lecture F.
